\section{Related Work}
\label{sec:related}

Tokenomic design has been studied extensively in the blockchain literature. Catalini and Gans \cite{catalini2016} provide foundational analysis of token-mediated markets. Cong et al.\ \cite{cong2021tokenomics} formalize the interaction between token prices and platform adoption. Sockin and Xiong \cite{sockin2023} analyze decentralized platform tokens as both currencies and securities.

Fee mechanism design was advanced significantly by Roughgarden's analysis of EIP-1559 \cite{roughgarden2020eip1559}, which proved incentive compatibility of the base fee adjustment. Chung and Shi \cite{chung2023} extended this to multi-dimensional fee markets relevant to multi-chain architectures like Lux.

Staking economics have been analyzed by Fanti et al.\ \cite{fanti2019economics}, who study compounding in Proof-of-Stake systems, and by Neuder et al.\ \cite{neuder2021}, who examine the interaction between staking yields and DeFi opportunities. Our adaptive reward function builds on these foundations while addressing multi-chain dynamics.

